\documentclass[11pt,a4paper]{article}
\usepackage[T1]{fontenc}\usepackage{lmodern,amsmath,amssymb,amsthm,mathtools,booktabs,microtype}
\usepackage[margin=27mm]{geometry}\usepackage[colorlinks=true,linkcolor=blue,urlcolor=blue,citecolor=blue]{hyperref}
\newtheorem{theorem}{Theorem}[section]\newtheorem{lemma}[theorem]{Lemma}\newtheorem{proposition}[theorem]{Proposition}
\newcommand{\Q}{\mathbb Q}\newcommand{\Z}{\mathbb Z}\DeclareMathOperator{\Tr}{Tr}\DeclareMathOperator{\Nm}{Nm}\DeclareMathOperator{\rad}{rad}
\setlength{\emergencystretch}{2em}
\title{An 81-variable quartic under rational $abc$}
\author{}\date{13 November 2026}
\begin{document}\maketitle
\begin{abstract}
Under rational $abc$ and the specified elliptic-group, prime-selection
and pole-realization premises, we construct a family of quartics in
81 total rational variables whose input projection is the integers
for every sufficiently large fixed odd coefficient shift. The family
is the sum of 73 primitive quadratic squares, has exactly 941
coordinate monomials, and has coefficient height
$\max(4153456,B^2+797551)$ for every fixed odd $B\ge3$.
The sufficient defining shift is ineffective.
The reduction from 87 variables uses two reciprocal elliptic points,
recovers their difference instead of storing a separate root point,
and permits half-integral indices. A proved trace-height theorem
covers these additional points. Additive comparison arguments remove
one more stored scale. Identity, tangent and torsion cases are retained
until the arithmetic argument treats them.
\end{abstract}

\section{Statement and exact premises}
All variable counts include the input. Heights of integer polynomials
are maximum absolute coefficients after coefficient parameters have
been substituted. Put
\[
 K=\Q(\rho),\quad\rho^2=2,\quad d=3-\rho,\quad D=d^2=11-6\rho,
 \qquad E:y^2=x^3-Dx,
\]
\[
 R_0=(1+\rho,2\rho),\qquad
 P=2R_0=(a,b)=\left(\frac{57-28\rho}{8},\frac{-660+417\rho}{32}\right).
\]
The retained arithmetic premises have the following scope.

\paragraph{Elliptic group.}
We assume $2E(K)=\Z P$ and $P$ non-torsion. Hence every point has
an index in $\frac12\Z$, as proved in \cite[Corollary 1.3]{TraceAll}.
The reciprocal pair below is not restricted to $2E(K)$.

\paragraph{Pole obstruction and realization.}
For $b_0\ne0$ let $B_0=17b_0$. Rational solutions of
\begin{equation}\label{eq:originalpole}
 p_0^2=B_0(1-e_0)(1+3e_0),\qquad
 B_0q_0^2=(s_0^2-B_0e_0(B_0-b_0^2))(s_0^2+3B_0e_0(B_0-b_0^2))
\end{equation}
exclude negative odd valuations of $b_0$ outside $2,3,17$.
We retain the converse with $p_0q_0s_0\ne0$ for prepared inputs
\begin{equation}\label{eq:prepared}
\begin{gathered}
0<b_0<17,\qquad b_0\text{ has square reduced denominator},\\
b_0\in\Q_2^{\times2}\cap\Q_3^{\times2}\cap\Q_{17}^{\times2},\qquad
v_2(b_0)\ge3,\quad v_3(b_0)\ge1,\quad v_{17}(b_0)\ge2.
\end{gathered}
\end{equation}
These are the same obstruction and prepared nonzero-realization
premises as the predecessor \cite{Old87}; their rank and descent
foundations are not independently re-proved here.

\paragraph{Prime selection.}
We retain the fixed-data simultaneous-prime statement of
\cite{Dependencies,Selection}. It applies to any fixed finite list
of distinct rational contact coefficients after a fixed square
clearing of denominators. It allows prescribed divisibility of the
square denominator, finite avoidance, an arbitrarily large fixed
lower bound on the rational argument, and the Gaussian congruences
and conjugate annihilating orders needed for the common index.
In the source normalization these orders $U_*,V_*$ and preparation
modulus $L_*$ have
\[
\gcd(U_*,V_*)=1,\quad\gcd(U_*,L_*)=16,\quad\gcd(V_*,L_*)=1,
\]
so $h\equiv0\pmod{U_*}$, $h\equiv1\pmod{V_*}$ and
$h\equiv h_*\pmod{L_*}$ are compatible when $64\mid h_*$.
The source also retains the prescribed compact real neighborhoods
in that progression. All finite data and rational coefficients are
fixed before prime selection. We check the new coefficient list and
all subsequent uses in Section~\ref{sec:complete}.

\paragraph{Additional hypothesis.}
Ordinary rational $abc$ is assumed. The eight-contact estimate below
is not asserted to follow from the separate weak-$abc$/Szpiro
hypothesis of the thirteen-variable family. No modular contact-height
premise is needed here.

\begin{theorem}\label{thm:main}
Under these premises there is a computable family
$P_B\in\Z[X,Y_1,\ldots,Y_{80}]$, indexed by odd integers $B\ge3$,
such that every member has degree four, content one, 941 nonzero
monomials, and height $\max(4153456,B^2+797551)$.
There is a fixed $B_*$ such that for every odd $B\ge B_*$,
\[
 x\in\Z\quad\Longleftrightarrow\quad
 \exists\boldsymbol y\in\Q^{80}:P_B(x,\boldsymbol y)=0.
\]
The coefficient parameter $B$ is not a quantified variable. No
numerical sufficient $B_*$ is given; the $B=3$ table is an algebraic
fixture, not a certified defining member.
\end{theorem}

\section{The quadratic system}
A field variable denotes two rational coordinates in the basis
$1,\rho$. A bar denotes conjugation. The notation $S^2$ denotes
an ordinary product, whereas the variables named $F_H,G_H,F_U,F_V$
are stored field values. Every displayed equation has degree at most
two in the retained rational coordinates.

\subsection{The common reference point}
Put $A_P=a^2+D$, $B_P=a^2-D$. Retain $S,T_H,f_H,g_H\in K$ and impose
\begin{align}
 2bD(aT_H-b)f_H&=aA_PST_H+bA_PS-bB_PT_H+4a^2D,\\
 2bD(aS+b)g_H&=-aA_PST_H-bB_PS+bA_PT_H-4a^2D,\\
 D^2f_Hg_H+aD(f_H+g_H)+aST_H+b(S-T_H)+D&=0.\label{eq:common}
\end{align}
These are the complete fixed-difference equations of \cite{Common}.
They give $H,H-P$ finite with nonzero abscissas and
\[
 S=y(H)/x(H),\quad T_H=y(H-P)/x(H-P),\quad
 f_H=1/x(H),\quad g_H=1/x(H-P).
\]
In particular $f_Hg_H\ne0$, but we do not assume $S T_H\ne0$ on
all components.

For explicitness, the exceptional-leading-coefficient check is as
follows. Put
\[
 \Phi=aS^2T_H^2-(a^2-D)(S^2+T_H^2)+2(a^2+D)ST_H+4aD,
\]
and define
\begin{align*}
2bDf_0&=-(a^2+D)S+(a^2-D)T_H-bS^2-aS^2T_H,\\
2bDg_0&=(a^2+D)T_H-(a^2-D)S-bT_H^2+aST_H^2.
\end{align*}
With $L_f,L_g$ the two leading linear coefficients, and $r_f,r_g$
the first two residuals, the last residual $R$ satisfies
\[
L_fL_gR=D^2r_fr_g+(aDL_g+D^2Q_g)r_f+(aDL_f+D^2Q_f)r_g
                 -aD^2C_P\Phi,
\]
where $Q_f,Q_g$ are the displayed right sides and
$C_P=D^2+6a^2D-3a^4\ne0$ (its norm is
$-8203357319351/16777216$). Hence $\Phi=0$ at every solution.
If $L_fL_g\ne0$, then $f_H=f_0,g_H=g_0$. If $L_f=0$, then
$T_H=b/a$, $L_g=-aDC_P/(a^2+D)\ne0$, $g_H=1/a$, and the last
equation has nonzero $f_H$ coefficient $D(a^2+D)/a$. It again gives
$f_H=f_0$. The other case is symmetric. The identities
\[
 Df_0^2+S^2f_0-1=\frac{a^2S^2-b^2}{4ab^2D}\Phi,
 \qquad
 Dg_0^2+T_H^2g_0-1=\frac{a^2T_H^2-b^2}{4ab^2D}\Phi
\]
exclude zero reciprocals. The chord identity
$a(f_0T_H+g_0S)+b(g_0-f_0)=S+T_H$, together with the abscissa
chord formula, gives $H-P$; when $x(H)=a$, the only allowed case
is $H=-P$, and the tangent formula gives $H-P=-2P$.
The converse follows from the same identities. Thus no open-subset
assumption is hidden in \eqref{eq:common}.

\subsection{Two points and their difference}
Retain $f,h,g,k,s\in K$, with stored squares $F_U,F_V$, and impose
\begin{equation}\label{eq:pair}
 F_U=f^2,\quad F_V=g^2,\quad h^2=f-DfF_U,\quad
 s(f-g)=h+k,\quad s(h-k)=1-D(F_U+fg+F_V).
\end{equation}
The complete algebraic proof in \cite[Proposition 2.1]{TraceAll}
identifies points $U,V\in E(K)$, including the identity, with
$U\ne V$. Their finite difference $A=U-V$ has
\[
 c=x(A)=s^2+D(f+g),\qquad y(A)=Dh-s(Df+c).
\]
Neither $U,V$ nor $A$ is assumed to have square abscissa. The theorem
includes the tangent case $f=g$, and excludes $(f,h)=(g,k)=(0,0)$.
Define one retained rational coordinate
\begin{equation}\label{eq:traceheightcoord}
                           T=2\operatorname{Re}(s^2+D(f+g))=\Tr x(A).
\end{equation}
This is a quadratic equation. The index $n$ of $A$ belongs to
$\frac12\Z$, and may be zero.

\subsection{Ratios, comparisons and the additive scale}
Retain $L,L_0\in K$ with
\begin{align}
 hL+fS&=0,& (1-DF_U)L+hS&=0,\\
 kL_0+gT_H&=0,& (1-DF_V)L_0+kT_H&=0.\label{eq:ratios}
\end{align}
At formal nonidentity points these are the usual ratios, and at an
identity numerator they force the ratio to be zero. No division by
$h$ or $k$ is made on torsion components.

Retain the input $X\in\Q$, a comparison coordinate $I_0\in\Q$ and a
slope $\lambda\in\Q$. Write $L=L_r+\rho L_i$ and impose
\[
 I_0-L_r+X^2+B=0,
 \qquad \beta=L_r-I_0=X^2+B,
 \qquad I=I_0+\rho L_i,
 \qquad C=L-\overline{L_0}.
\]
Retain $F_H,G_H,C_2,I_2\in K$ and two scalar height coordinates
$u_h,v_h$, with
\[
 F_H=f_H^2,\quad G_H=g_H^2,\quad C_2=C^2,\quad I_2=I^2,
 \qquad u_h=\beta^2+T^2,\quad v_h=\beta T.
\]
The scale $a_1=\beta^4+T^4=u_h^2-2v_h^2$ is not retained as a
coordinate. Use arguments
\begin{equation}\label{eq:arguments}
            t_0=\lambda,\qquad t_1-\lambda=u_h^2-2v_h^2.
\end{equation}
Only $t_1$ is a further coordinate. Notice that $a_1\ge B^4\ge81$.

\subsection{Contact and pole equations}
For each argument $t=t_j$ retain scalars $R,Z,w,w_2,w_3$, field
values $A_1,A_2,\Omega=\omega_0+\rho\omega_1$, and scalars
$\eta,b_0,e,u,p,v,q,p_2,q_2$. Impose
\begin{align}
 R&=t^2-2t,\qquad RZ+Z-R^2+35R=0,\qquad (R+1)w=Z^2+24Z,\\
 w_2&=w^2,\qquad w_3=ww_2,\\
 w_3A_1&=w_2^2+w_3C_2+w_2F_H+F_U^2,\\
 w_3A_2&=w_2^2+w_3I_2+w_2\overline{G_H}+\overline{F_V}^{\,2},\\
 (R+1)\Omega&=A_1A_2,\\
 \omega_0^2-2\omega_0\eta-2\omega_1^2&=0,\qquad w_3b_0=\eta^2,\\
 p_2&=p^2,\qquad 17b_0p_2=289+34e-3e^2,\\
 e(17-b_0)&=17uv,\qquad q_2=q^2,\\
 17b_0q_2&=u^2+34uv-867v^2.\label{eq:blocks}
\end{align}
In each factor equation a scalar is interpreted as a field element.
The two blocks have separate pole coordinates.
For $b_0\ne0$ the last five equations imply \eqref{eq:originalpole}
with
\[
 e_0=e/17,\quad p_0=b_0p,\quad s_0=b_0u,\quad q_0=b_0^2u q.
\]
Conversely a prepared witness with $s_0\ne0$ gives these coordinates,
using $v=e(17-b_0)/(17u)$. Zero values on other components are not
silently divided out.

\section{Domains and local extraction}
Put $r=R+1$, $z=Z+37$. The contact equations give
\[
 r=(t-1)^2,\quad r(z-r)=36,\quad rw=z^2-50z+481.
\]
Thus $r\ne0$. With
\[
 P_4(r)=(r-1)(r-4)(r-9)(r-36),\qquad
 \Theta=\{-5,-2,-1,0,2,3,4,7\},
\]
one has $P_4(r)/r^2=z^2-50z+481$ and
$F(t):=\prod_{\theta\in\Theta}(t-\theta)=wr^3$.
If $w=0$, the factor equations give $F_U^2=\overline{F_V}^{\,2}=0$,
hence $f=g=0$; the cubic gives $h=0$, the secant gives $k=0$, and
the last equation of \eqref{eq:pair} gives $0=1$. Therefore
$r w F(t)\ne0$ at every rational solution, and $t\notin\Theta$.

For interpretation only, define
\[
 q_* =w^{-1}=\frac{(t-1)^6}{F(t)},\quad
 R_* =\frac{F(t)}{(t-1)^{10}},\quad
 T_i=A_i/w,\quad W=T_1T_2=\frac r{w^2}\Omega,
 \quad\tau=\frac r{w^2}\eta.
\]
Then
\begin{align*}
T_1&=1+q_*C^2+q_*^2 f_H^2+q_*^4f^4,\\
T_2&=1+q_*I^2+q_*^2\overline{g_H}^{\,2}+q_*^4\overline g^{\,4},
\qquad b_0=R_*\tau^2.
\end{align*}
The cleared trace equation is preserved under this common rescaling.
None of $W,\eta,b_0$ is assumed globally nonzero at this stage.

For $t\notin\Theta$, let $M(t)$ be the product of primes $p>76176$
with positive odd valuation at one of $t-\theta$. At such a prime,
all contact differences and $t-1$ are units. Write
$\nu=v_p(F(t))>0$, odd; then $v_p(q_*)=-\nu$.
Every nonzero reciprocal abscissa has even valuation at these good
places. This follows directly from $h^2=f-Df^3$ if $v(f)\ne0$;
valuation zero is even too. The same statement holds for the common
reference point. A zero reciprocal contributes no nonconstant term.

The finite nonconstant valuations in each factor are respectively
odd, $2$ modulo $4$, and $4$ modulo $8$, and differ from the
constant term of valuation zero. Thus both factors are nonzero and
have valuation at most zero at every place above $p$. Let
$\ell=\max_{\mathfrak p\mid p}v_{\mathfrak p}(W)\le0$.
The trace equation forces the rational component of $W$ and $\tau$
to be nonzero, and
\[
 \tau=\Tr(W^{-1})^{-1},\qquad
 v_p(b_0)=\nu-2v_p\Tr(W^{-1})\le\nu+2\ell.
\]
This valuation is odd. The pole obstruction forces it positive, so
$\ell>-\nu/2$. At a place attaining $\ell$, each nonconstant
summand valuation exceeds $-\nu/2$. Consequently
\[
 v(C),v(I)>\nu/4,\quad
 v(f_H),v(\overline{g_H})>3\nu/4,\quad
 v(f),v(\overline g)>7\nu/8.
\]
The reciprocal valuations are positive even integers, and $C,I$
have order at least one. In particular the four indicated points are
in the appropriate primary or conjugate formal groups, with the
identity convention at zero reciprocals.

Write the indices of $H,U,V,A$ as $h_*,u_*,u_*-n,n$ in
$\frac12\Z$. At a consumed prime, $H,H-P$ are nonidentity formal
points and therefore non-torsion. Multiplying the integer-index
group identities by two kills every possible two-torsion component.
The full formal-logarithm ratios are
\[
                u_*/h_*,\qquad (u_*-n)/(h_*-1).
\]
At least one denominator is a $p$-adic unit, and the numerator
indices are $p$-integral. Thus one full ratio is integral.
The order-four relative logarithm estimate in \cite{TraceAll}
first makes its plain ratio integral. The test $C\equiv0\pmod p$
then gives integrality of the other plain ratio, hence of its full
ratio. Only after this step are the relative errors used as absolute
congruences. Subtracting gives
\[
 n=h_*\frac{u_*}{h_*}-(h_*-1)\frac{u_*-n}{h_*-1}
          \equiv L\equiv\beta\pmod p.
\]
This proves that, if $\beta\ne n$ and $N$ is the positive reduced
numerator of $\beta-n$, then
\begin{equation}\label{eq:div}
                         M(t_0)\mid N,\qquad M(t_1)\mid N.
\end{equation}
Inert primes are included: if the required formal groups are
incompatible, no solution occurs at that prime. No assertion that
the two radicals are coprime is needed.

\section{Heights and the uniform sufficient shift}
\begin{lemma}\label{lem:abc}
Under rational $abc$, fixed constants $C_0,C_1$ satisfy
\[
 h(t)\le\tfrac13\log M(t)+C_0\quad(t\notin\Theta),\qquad
 h(1:u:v)\le\tfrac43 h(u^4+v^4)+C_1\quad(u,v\in\Q).
\]
\end{lemma}
\begin{proof}
For $t=a/b$ reduced use $G(a,b)=\prod_{\theta\in\Theta}(a-\theta b)$.
Granville's Proposition 1~\cite{Granville}, with exponent loss $1/2$,
gives $\rad(G)\gg\max(|a|,b)^{11/2}$. Since
$\rad(m)^2\le |m|\operatorname{sf}(m)$ and
$|G|\ll\max(|a|,b)^8$, its odd part is at least a fixed multiple of
$\max(|a|,b)^3$. Outside the fixed exceptional primes it is exactly
$M(t)$: two primitive contact factors cannot vanish together, and a
prime in $b$ divides none. This proves the first assertion, including
bounded heights by increasing $C_0$. The form has eight distinct
factors and $G(a,b)\ne0$, so all source hypotheses are met.
The author's corrigendum changes Theorems 8.1 and 8.3, not this
proposition; neither corrected theorem is used here.

For the second assertion, the zero-coordinate cases are immediate.
Otherwise write $u=gA/D_0$, $v=gB_0/D_0$ with
$\gcd(A,B_0)=\gcd(g,D_0)=1$, $g,D_0>0$. Set
$H=\max(|A|,|B_0|)$, $N_0=A^4+B_0^4$, and
$k=\gcd(N_0,D_0^4)$, $m_0=N_0/k$, $n_0=D_0^4/k$.
Then $u^4+v^4=g^4m_0/n_0$ is reduced. Put
$M=\max(g^4m_0,n_0)$. Rational $abc$ at exponent $16/15$ and
$\rad(N_0)\le D_0m_0$ give
$H^{7/4}\le K^{15/16}D_0m_0$. Also
$D_0^4\le2n_0H^4/m_0$, whence
$H^{3/4}\le2^{1/4}K^{15/16}m_0^{3/4}n_0^{1/4}$.
Finally $\max(gH,D_0)\le2^{1/4}H(M/m_0)^{1/4}$. Taking logarithms
bounds $h(1:u:v)$ by a fixed constant plus
$\frac34\log m_0+\frac13\log n_0+\frac14\log M\le\frac43\log M$.
This accounts for all common numerator and denominator cancellation.
\end{proof}

Suppose $\beta\ne n$. By \eqref{eq:arguments}, \eqref{eq:div} and
the elementary additive height inequality,
\[
 h(a_1)\le h(t_0)+h(t_1)+\log2\le\tfrac23\log N+2C_0+\log2.
\]
Writing $H_2=h(1:\beta:T)$, the preceding lemma yields
\[
                         H_2\le\tfrac89\log N+C_2.
\]
Because $2n\in\Z$, the reduced numerator satisfies
$N\le2(1+|n|)e^{h(\beta)}$. Hence
\[
             H_2\le8\log(2(1+|n|))+9C_2.
\]
The finite point $A=U-V$ satisfies the all-point trace theorem
\cite[Corollary 1.3]{TraceAll}, including torsion and half-integral
indices. Thus
\[
 n^2\le8A_0\log(2(1+|n|))+9A_0C_2+J_0.
\]
This bounds $|n|$ by a fixed constant and then bounds $H_2$ by a fixed
constant, independent of $B$ and every witness. But
$h(\beta)\ge\log\beta\ge\log B$, since $\beta=X^2+B$.
One sufficiently large fixed $B_*$ excludes $\beta\ne n$ for all
$B\ge B_*$. Therefore $\beta\in\frac12\Z$.
Writing $X=a/b$ reduced gives $b^2\mid2a^2$, hence $b^2\mid2$ and
$b=1$. This proves soundness. The replacement of $2X^2+B$ by
$X^2+B$ is essential for this denominator argument with half-indices.

\section{Completeness}\label{sec:complete}
Fix an integer $X$ and an odd $B\ge3$. Set $n=\beta=X^2+B>0$ and
$A=[n]P$. This fixes $T=\Tr x(A)$ and $a_1=\beta^4+T^4\ge81$
before selecting primes. Choose a nonexceptional reference multiple
$H=[h]P$, with its initial index divisible by $64$, in fixed compact
real and finite-place neighborhoods. Set $U=[hn]P$ and
$V=[(h-1)n]P$. All are in the doubled group, so the usual
square-abscissa assertions are available for these particular
witnesses, although they were not imposed on all solutions.
Avoid finitely many forbidden reference points and denominators.
The reciprocal-pair secant or tangent lemma then supplies its
coordinates; $A$ is never the identity.

The new fixed list of contact coefficients in the slope variable is
\[
                    \Theta\ \cup\ (\Theta-a_1).
\]
The two lists are disjoint since their span is $12$ and $a_1\ge81$.
Choose an integer $D_1>0$ whose square clears these rational
coefficients, sufficiently divisible at every preparation prime.
Apply the retained prime-selection input to this list with
$\lambda=u/(D_1^2w^2)$. Its prime numerators are exactly
\[
 u-\theta D_1^2w^2,\qquad
 u+(a_1-\theta)D_1^2w^2\qquad(\theta\in\Theta).
\]
The zero contact in the first list makes $u$ one selected prime.
The finite avoidance data include all denominators and contact
collisions. The prescribed conjugate reduction orders and the
compatible common-index progression make $H,U$ formal at the
primary places, and $H-P,V$ formal at the opposite places. Both
full ratios equal $n$, so the reciprocal tests have depth at least
two and $C,I$ have depth at least four. With $v(q_*)=-1$, every
nonconstant factor term is positive at the selected place. At the
opposite place each factor is either a unit of residue one or has
negative valuation. Thus $\Tr(W^{-1})$ reduces to one or two and is
a unit. The actual pole input has valuation one there.

Put $v_0=D_1^2w^2$. For $a\in\{0,a_1\}$ the exact identity is
\[
 R_*(\lambda+a)=
 \frac{\prod_{\theta\in\Theta}(u+(a-\theta)v_0)}
 {v_0^8(\lambda+a-1)^{10}}.
\]
All factors other than the selected prime numerators are rational
squares. The trace changes every valuation by an even number.
Consequently the pole input has square reduced denominator.
At $2,3,17$ choose denominator divisibility so the argument dominates
the fixed $a_1$ $p$-adically. Then $q_*,R_*\sim t^{-2}$, all weighted
tests are small, and $\Tr(W^{-1})$ is arbitrarily close to two.
This gives the required local squareclasses and depths in
\eqref{eq:prepared}. On the compact real neighborhoods, choosing an
arbitrarily large fixed lower bound on $\lambda$ makes both arguments
large and positive; $W\to1$ and $b_0\sim1/(4t^2)$, so
$0<b_0<17$. These bounds are chosen before the prime selection.
Independent prepared pole witnesses now fill both blocks.

No parity restriction on $n$ was used: multiplication by an even or
odd integer preserves the formal group, and the full ratios are the
same integer $n$. The construction works for $X=0$ as well. It proves
completeness for every integer input and every allowed fixed shift.
Together with the uniform soundness threshold it proves
Theorem~\ref{thm:main}.

\section{Counts, coefficients and verification}
The retained rational-coordinate census is
\begin{center}\begin{tabular}{lr}\toprule
Input, real comparison and slope&3\\
Common reference chart&8\\
Reciprocal pair and secant parameter&10\\
Two field ratios&4\\
Four reciprocal squares and two comparison squares&12\\
Trace and two height outputs&3\\
Two contact/pole blocks&40\\
Second argument&1\\\midrule Total&81\\\bottomrule
\end{tabular}\end{center}
There are 73 rational quadratic equations. Clear each to a primitive
integer row with positive first coefficient in the specified order.
For the first six common-chart rows apply the fixed invertible
rational matrix in \texttt{common-row-reduction.json}, whose
determinant is $-1/362963666515968$, and normalize again. This changes
no zero set. The sum of the 73 row squares is $P_B$.

The directory \texttt{computation/quartic/abc81\_v25} contains the
complete coordinate and row table and its fully expanded quartic.
Exact collection gives 942 parameter-separated terms and 941
coordinate monomials for every positive $B$. The coefficient of
$X^4$ is one, proving exact degree and content. The constant
coefficient is $B^2+797551$; the largest fixed coefficient is
$4153456$. Every remaining parameter coefficient is a linear
multiple of $B$, bounded by the maximum of these two values.
Thus the stated height formula is exact. For odd $B\ge1833$ it
is $B^2+797551$; this is only a coefficient threshold.

A separate implementation reconstructs all 73 equations directly
from the field formulas using symbolic polynomial arithmetic, checks
the rational row matrix, and independently expands every quartic
coefficient. It checks the universal point and contact identities
and rejects a changed coefficient, missing row and corrupted quartic.
The 702 exact elliptic examples described in \cite{TraceAll} include
the exceptional branches used here. The computations do not establish
$abc$, the inherited arithmetic premises, a numerical defining shift,
or non-author acceptance of the proof.

\begin{thebibliography}{9}
\bibitem{TraceAll} Companion paper, \emph{Trace heights on quadratic
twists with $j=1728$}, 13 November 2026,
\nolinkurl{papers/trace_heights/trace_on_all_points_20261113_v25.tex}.
\bibitem{Granville} Andrew Granville, \emph{Rational and integral points
on quadratic twists of a given hyperelliptic curve}, International
Mathematics Research Notices (2007), rnm027, Proposition 1; author's
version, printed p.~6, \url{https://dms.umontreal.ca/~andrew/PDF/hyperelliptics.pdf}.
See also \url{https://dms.umontreal.ca/~andrew/PDF/hyperellCorr.pdf};
its corrections to Theorems 8.1 and 8.3 are not used here.
\bibitem{Old87} Companion paper, \emph{An 87-variable quartic under
$abc$}, \nolinkurl{papers/quartic_existence/quartic_abc87_20261110_v21.tex},
repository commit \texttt{fa09c9c6c97982c13265e7088d95cdb61b3e90cf}.
\bibitem{Common} Companion paper, \emph{Fixed-difference elliptic
charts}, \nolinkurl{papers/elliptic_charts/fixed_difference_charts_20261110_v18.tex}.
\bibitem{Dependencies} Companion paper,
\nolinkurl{papers/arithmetic_dependencies/arithmetic_dependencies_20261108_v8.tex}.
\bibitem{Selection} Companion paper,
\nolinkurl{papers/prime_selection/prime_selection_20261109_v12.tex}.
\end{thebibliography}\end{document}
